Dans un bécher ouvert posé sur une balance, on verse du vinaigre puis du bicarbonate de soude. La masse totale avant la réaction est de \qty{\v{m1}}{g}. Après l'effervescence, la balance indique \qty{\fpeval{\v{m1}-\v{dm}}}{g}. La réaction est :
\[ \text{NaHCO}_3 + \text{CH}_3\text{COOH} \rightarrow \text{NaCH}_3\text{COO} + \text{H}_2\text{O} + \text{CO}_2 \]
Masses molaires : $M(\text{CO}_2) = \qty{44.01}{g/mol}$, $M(\text{NaHCO}_3) = \qty{84.01}{g/mol}$.

\begin{enumerate}[label=\alph*)]
    \item Quelle masse de dioxyde de carbone s'est échappée ? Expliquez pourquoi la masse mesurée diminue, même si la loi de conservation de la masse est respectée.
    \item Combien de moles de dioxyde de carbone se sont formées ?
    \item Quelle masse de bicarbonate de soude a réagi ?
\end{enumerate}
